% ==============================================================================
% preface.tex — คำนำ และโครงสร้างหลักสูตรการฝึกอบรมเชิงปฏิบัติการ
% ==============================================================================
\chapter*{คำนำ}
\addcontentsline{toc}{chapter}{คำนำ}

การก้าวเข้าสู่ยุคปัญญาประดิษฐ์เชิงสร้างสรรค์ (Generative AI) และเทคโนโลยีปฏิสัมพันธ์เชิงพื้นที่ (Spatial Computing) ได้สร้างจุดเปลี่ยนครั้งประวัติศาสตร์ต่อการจัดการศึกษาวิทยาศาสตร์และเทคโนโลยี การจัดการเรียนรู้ในรูปแบบเดิมที่เน้นการบรรยายสองมิติผ่านกระดานหรือแผ่นสไลด์ ไม่สามารถตอบสนองต่อการทำความเข้าใจในปรากฏการณ์ทางกายภาพที่มีความซับซ้อน เช่น การไหลของกระแสไฟฟ้าในโครงข่ายตัวต้านทานหลายมิติ การเคลื่อนที่แบบฮาร์มอนิกอย่างง่าย หรือการประกอบวงจรทดลองที่มีความละเอียดอ่อนสูงได้ดีเพียงพอ

คู่มือการอบรมเชิงปฏิบัติการเรื่อง \textbf{“การออกแบบและพัฒนาสื่อการสอนสามมิติเสมือนจริงและความจริงเสริมด้วยปัญญาประดิษฐ์” (Designing and Developing 3D Virtual \& Augmented Reality Instructional Media with AI)} เล่มนี้ จัดทำขึ้นเพื่อเป็นพิมพ์เขียวเชิงเทคนิค (Technical Blueprint) และคู่มือภาคปฏิบัติการสำหรับคณาจารย์ ครูผู้สอนวิทยาศาสตร์ และนักพัฒนานวัตกรรมการศึกษา โดยสังเคราะห์จากองค์ความรู้และการพัฒนาจริงในระบบ \textbf{RBRU Next Gen Lab 2026} ของศูนย์ความเป็นเลิศด้านสิ่งแวดล้อมและพลังงาน (LEQs) คณะวิทยาศาสตร์และเทคโนโลยี มหาวิทยาลัยราชภัฏรำไพพรรณี 

เนื้อหาภายในเล่มได้รับการออกแบบอย่างเป็นระบบตามแนวคิด **Active Learning Cycle** และ **Outcome-Based Education (OBE)** แบ่งออกเป็น 6 บทเรียนหลัก ครอบคลุมตั้งแต่การใช้ปัญญาประดิษฐ์ช่วยวิเคราะห์มาตรฐานตัวชี้วัดหลักสูตร, การเรนเดอร์ภาพสามมิติความสมจริงเชิงฟิสิกส์ (Three.js PBR Shaders), การสร้างเอนจินลดรูปวงจรไฟฟ้าแบบเวกเตอร์ SVG ละเอียดทุกขั้นตอน (Step-by-Step Circuit Reduction อิงมาตรฐานข้อสอบ 20th IJSO), การตรวจจับท่าทางมือไร้สัมผัส (Touchless AR MediaPipe Hands), จนถึงการติดตั้งเผยแพร่บนคลาวด์ CDN ระดับโลก

ผู้เขียนหวังเป็นอย่างยิ่งว่า คู่มือเล่มนี้จะช่วยเสริมสร้างขีดความสามารถของครูและอาจารย์ ให้สามารถแปลงความรู้เชิงทฤษฎีทางวิทยาศาสตร์ ให้กลายเป็นสื่อดิจิทัลสามมิติเสมือนจริงที่นักเรียนสามารถสัมผัส ทดลอง และมีปฏิสัมพันธ์ได้อย่างไร้ขีดจำกัด อันเป็นรากฐานสำคัญในการขับเคลื่อนคุณภาพการศึกษาไทยสู่ระดับสากล

\vspace{1.5cm}
\begin{flushright}
\textbf{ผู้ช่วยศาสตราจารย์ ดร.ชีวะ ทัศนา}\\
ศูนย์ความเป็นเลิศด้านสิ่งแวดล้อมและพลังงาน (LEQs)\\
คณะวิทยาศาสตร์และเทคโนโลยี มหาวิทยาลัยราชภัฏรำไพพรรณี\\
ตุลาคม พ.ศ. 2569
\end{flushright}

\newpage

% ==============================================================================
% ตารางกำหนดการอบรมเชิงปฏิบัติการ 2 วัน (12 ชั่วโมง)
% ==============================================================================
\chapter*{โครงสร้างหลักสูตรการอบรมเชิงปฏิบัติการ (12 ชั่วโมง)}
\addcontentsline{toc}{chapter}{โครงสร้างหลักสูตรการอบรมเชิงปฏิบัติการ (12 ชั่วโมง)}

หลักสูตรการอบรมเชิงปฏิบัติการนี้จัดขึ้นเป็นเวลา 2 วัน รวม 12 ชั่วโมงการเรียนรู้เชิงลึก (Intensive Hands-on Workshop) มุ่งเน้นการลงมือปฏิบัติจริงแบบ 100\% (BYOD: Bring Your Own Device) โดยผู้เข้ารับการอบรมจะได้รับเว็บแอปพลิเคชันห้องปฏิบัติการเสมือนจริงที่ทำงานได้สมบูรณ์เป็นของตนเองเมื่อสิ้นสุดการอบรม

\vspace{0.3cm}

\begin{table}[h!]
\centering
\small
\renewcommand{\arraystretch}{1.25}
\begin{tabular}{p{2.2cm}p{4.2cm}p{6.0cm}}
\toprule
\textbf{ช่วงเวลา} & \textbf{หัวข้อการเรียนรู้} & \textbf{กิจกรรมและผลผลิตเชิงประจักษ์} \\
\midrule
\multicolumn{3}{l}{\textbf{\color{rbruNavy}วันแรก: สถาปัตยกรรม 3D Virtual Lab และการร่างคำสั่ง AI (Day 1)}} \\
09:00 -- 10:30 น. & บทที่ 1: วิวัฒนาการและทฤษฎีสื่อ 3 มิติเสมือนจริง & วิเคราะห์ข้อจำกัดสื่อ 2D สู่มิติ 3D เสมือนจริง และ Cognitive Load Theory \\
10:45 -- 12:00 น. & บทที่ 2: AI Prompting เพื่อออกแบบสื่อการสอน & การใช้ CLEAR Prompting และ Bloom's Taxonomy เพื่อสังเคราะห์โค้ดสื่อการสอน \\
13:00 -- 14:45 น. & บทที่ 3: สถาปัตยกรรม Three.js PBR Shaders & สร้างฉากทดลอง 3D จัดแสงเงาตกกระทบสมจริง และการสร้างอุปกรณ์ฟิสิกส์ \\
15:00 -- 16:30 น. & ปฏิบัติการ: ระบบลากย้ายอุปกรณ์อิสระ & การทำ Raycasting บนระนาบ $y=0$ และสายไฟยืดหยุ่น Catmull-Rom Leads \\
\midrule
\multicolumn{3}{l}{\textbf{\color{rbruNavy}วันที่สอง: วงจรโอลิมปิก, ระบบ AR ไร้สัมผัส และการเผยแพร่ (Day 2)}} \\
09:00 -- 10:30 น. & บทที่ 4: เอนจินลดรูปวงจรไฟฟ้า SVG พลวัต & คลี่วงแหวน 6 เหลี่ยม IJSO ออกเป็น 2 กิ่ง และการคำนวณเศษส่วนแท้ Euclidean GCD \\
10:45 -- 12:00 น. & ปฏิบัติการ: แถบสีตัวต้านทานและหน้าปัดเข็ม & สร้างโค้ด EIA 4 แถบสีพลวัต และฟิสิกส์สปริงหน่วงเข็มมิเตอร์ (Second-Order) \\
13:00 -- 14:45 น. & บทที่ 5: การตรวจจับมือไร้สัมผัส AR MediaPipe & เชื่อมต่อกล้อง Webcam ตรวจจับ 21 จุดกระดูกมือ และระบบจีบนิ้ว Pinch Drag \\
15:00 -- 16:30 น. & บทที่ 6: การติดตั้ง CDN และประเมินผลสัมฤทธิ์ & Deploy บน GitHub Pages Global CDN ตรวจสอบ 0 Errors และวัดผลการเรียนรู้ \\
\bottomrule
\end{tabular}
\end{table}

\cleardoublepage
